\section{Magnetic domains and joint-prime gluing}\label{sec:gluing}
\subsection{What higher magnetic powers would require}
The domain statement for one insertion is sufficient for
Theorem~\ref{thm:schur-prime}. It does not place the prepared state in
the domain of every negative magnetic power.
\begin{proposition}[Formal higher-power coefficient criterion]
\label{prop:magnetic-powers}
For $k\geq1$, iteration of the coefficient action gives the formal series
\begin{equation}\label{eq:magnetic-k}
 (u_-^k\eta_a)_{-k}(\zeta)=\sqrt{1-a^2}
 \frac{\prod_{j=0}^{k-1}(1+q^{k-1-2j}\zeta)}{1+aq^{-k}\zeta}.
\end{equation}
Its Taylor coefficient sequence belongs to $\ell^2$ if and only if
\begin{equation}\label{eq:magnetic-criterion}
 a<q^k\quad\hbox{or}\quad a\in\{q,q^3,\ldots,q^{2k-1}\}.
\end{equation}
For fixed $0<a<1$, square summability of all these formal outputs is
equivalent to $a=q^{2j+1}$ for some integer $j\geq0$.
\end{proposition}
\begin{proof}
Induction in \eqref{eq:schur-actions} gives \eqref{eq:magnetic-k}.
For every coefficient of degree $n\geq k$, expanding the denominator
shows that the coefficient equals
\begin{equation}\label{eq:magnetic-tail}
 \sqrt{1-a^2}(-aq^{-k})^n
 \prod_{j=0}^{k-1}(1-q^{2j+1}/a).
\end{equation}
If the product is nonzero, the geometric tail is square summable
exactly when $aq^{-k}<1$. If the product vanishes, the numerator
cancels the denominator and the output is a polynomial. This proves
\eqref{eq:magnetic-criterion}. For all $k$, the inequality eventually
fails at any fixed $a>0$, so an exceptional value is necessary.
If $a=q^{2j+1}$, it satisfies the strict inequality for $k\leq j$
and the cancellation condition for $k\geq j+1$, proving sufficiency
for the formal coefficient statement.
\end{proof}

The exceptional square-summable outputs do not by themselves prove
membership in powers of a chosen closed operator extension. A divergent
coefficient tail is, however, a necessary domain obstruction whenever
the prescribed action on Fourier coefficients is retained. This
distinction prevents meromorphic continuation across a pole from being
mistaken for an operator-domain argument.

For the prime states $a_p=qp^{-1/2}$, requiring all negative powers
would force $p^{-1/2}=q^{2j_p}$. This cannot hold simultaneously for
$p=2$ and $p=3$, since it would imply a nontrivial equality between
an integer power of $2$ and an integer power of $3$. It does not affect
the single insertion already proved. More generally, for any fixed
finite $K>1$, choosing $q^{K-1}>2^{-1/2}$ puts every prime state in
the strictly decaying regime for all powers through $K$. In that
regime the Taylor-approximation argument of
Lemma~\ref{lem:one-magnetic} extends through the finite sequence of
shifts. A proposed junction must specify the finite words or actual
operator domains it uses; a field theory need not put every state in
the domain of every unbounded power.

\subsection{A two-prime obstruction for fixed output vectors}
For each prime write $r_p=p^{-1/2}$, $d_p=\log p$, and
\begin{equation}\label{eq:prime-filter}
 A_p(z)=\frac{1-z}{1-r_pz}
 =1-(1-r_p)\sum_{n\geq1}r_p^{n-1}z^n,
 \qquad c_p=\frac{d_pr_p(1+r_p)}{1-r_p}.
\end{equation}
The norm of $\sqrt{c_p}A_p$ is the multiplier
\eqref{eq:prime-weight}. Consider a proposed coherent sum with fixed
unit output vectors $v_2,v_3$:
\begin{equation}\label{eq:constant-vector-source}
 F\longmapsto\widehat F(\tau)
 \left[\sqrt{c_2}A_2(e^{id_2\tau})v_2
       +\sqrt{c_3}A_3(e^{id_3\tau})v_3\right],
\end{equation}
with output measure $\dd\tau/(2\pi)$.

\begin{proposition}[Mixed atom from constant-vector gluing]\label{prop:mixed}
In \eqref{eq:constant-vector-source}, the coefficient of
$e^{i(\log3-\log2)\tau}$ in the full norm multiplier is
\begin{equation}\label{eq:mixed-atom}
 \sqrt{c_2c_3}(1-r_2)(1-r_3)\ip{v_2}{v_3}.
\end{equation}
Consequently this source, together with independently added gamma,
contact and pole terms, cannot equal the two-prime Weil form on
$I_{5/4}$ unless $v_2\perp v_3$. In the orthogonal case its prime
part is exactly the sum of the separate positive references.
\end{proposition}
\begin{proof}
The first cross term in the squared norm is
$\sqrt{c_2c_3}\overline{A_2(e^{id_2\tau})}
A_3(e^{id_3\tau})\ip{v_2}{v_3}$.
Its first nonconstant coefficient from each factor gives
\eqref{eq:mixed-atom}. A second pair of nonnegative indices with
$n\log3-m\log2=\log3-\log2$ would give
$3^{n-1}=2^{m-1}$ and hence $n=m=1$. The reverse cross term cannot
give this positive frequency with nonnegative indices, and the
individual squared norms contain only integer multiples of $d_2$
or $d_3$.

The Fourier coefficient series of each factor is absolutely summable;
the cross kernel is therefore a finite-variation atomic measure. Other
mixed lengths may accumulate, but its atom at $\log(3/2)$ has the
specified nonzero mass if the overlap is nonzero. The Weil target has
no prime-power atom there. Its gamma and pole kernels are smooth near
this nonzero separation, and its contact is supported at zero, so none
can cancel the atom. Since $0<\log(3/2)<5/4$, restriction to $I_{5/4}$
still detects it. Equality of all polarized test-function pairings
would imply equality of their distribution kernels, a contradiction.
A purely imaginary overlap gives an imaginary antisymmetric kernel;
complex tests still detect it. Thus the overlap must vanish.
\end{proof}

The proposition does not apply to arbitrary coherent gluing. A joint
operator-valued junction may have additional terms which cancel mixed
returns as an identity. Producing and evaluating such a junction is a
new source construction. The failed fixed-vector ansatz gives no reason
to estimate its unwanted atoms as a small residual.

\subsection{The full comparison, with no omitted contact or pole}
Let $\mathcal P$ be any finite set containing all primes $p<e^L$.
Equations~\eqref{eq:target} and \eqref{eq:lifted-norm} give the exact
identity
\begin{align}
 Q_L[f]={}&K[E_Lf]+\sum_{p\in\mathcal P}B_{r_p,d_p}[E_Lf]
 \notag\\
 &+\left(w_0-\sum_{p\in\mathcal P}\kappa_{r_p,d_p}\right)\norm f^2
 +P_L[f].
 \label{eq:full-comparison}
\end{align}
Each term in the first line is independently positive. The second
line has not been realized as part of their joint norm. Extra primes
in $\mathcal P$ contribute only contacts on this interval, which
cancel algebraically against the displayed subtraction. This does
not make the subtraction an independently positive operation.

Thus the dressed result arrives at the same contact and pole problem
as the conservative prime reference. It supplies a concrete Schur
preparation and removes the earlier undressed-state mismatch, but
establishes no larger positivity interval. Unitary repackaging of the
orthogonal gamma and prime outputs leaves \eqref{eq:full-comparison}
unchanged. The next calculation must determine contact, poles and
prime interference within one complete positive preparation.
